\documentclass[11pt]{article}
\usepackage{amsmath,amssymb,amsthm,geometry,hyperref}
\geometry{margin=1in}
\newtheorem{claim}{Claim}
\title{The inverse Hall--Littlewood $P\to m$ matrix is \emph{not} in the
Wheeler--Zinn-Justin Hall--Littlewood papers\\[2pt]
\large and the proposed DFK locator is edition-confused}
\author{C.\ Vega (Clio)}
\date{2026-10-05}
\begin{document}
\maketitle

\noindent\fbox{\parbox{0.96\textwidth}{%
\textbf{From:} C.\ Vega (Clio)\quad\textbf{To:} R.\ (\texttt{grandpa-rick})\quad
\textbf{Date:} 2026-10-05\\
\textbf{Re:} your email of 2026-10-04 00:16, \emph{``s=1 block valuation law; is the
inverse HL $P\to m$ valuation known?''}\\
\textbf{Source commit:} \texttt{clio-vega/work-in-progress@3f06cec}\\
\textbf{Grade:} a \emph{negative} literature finding, measured at first hand on two
papers I hold, with controls. It is not a theorem.}}

\section{Your question (a), answered: no}

You asked whether the $t=0$ case of your block valuation law --- the $(1-q)$-adic
valuation of $[h_\mu]Q'_\lambda(x;q)$, an entry of the inverse Hall--Littlewood
$P\to m$ matrix, being $\ell(\lambda)-\kappa(\lambda,\mu)$ --- appears in the
Wheeler--Zinn-Justin Hall--Littlewood papers on my shelf.

\begin{claim}
It is not in either of them, and the reason is structural rather than accidental:
neither paper uses the monomial basis at all.
\end{claim}

I read both at first hand today. They are:
\begin{itemize}
\item Wheeler--Zinn-Justin, \emph{Hall polynomials, inverse Kostka polynomials and
puzzles}, \texttt{arXiv:1603.01815}; journal edition \emph{J.\ Combin.\ Theory Ser.\ A}
\textbf{159} (2018) 107--163. (I give both editions deliberately; see \S4.)
\item Zinn-Justin, \emph{Honeycombs for Hall polynomials}, \texttt{arXiv:1909.10720}v2.
\end{itemize}

\paragraph{What they actually compute.} WZJ study two one-parameter generalisations of
$c^\lambda_{\mu\nu}$, and \emph{both} live in bases other than $m$:
\begin{align}
P_\mu P_\nu &= \sum_\lambda f^\lambda_{\mu\nu}(t)\,P_\lambda
&&\text{(Hall polynomials; the $P$ basis),}\\
s_\mu P_\nu &= \sum_\lambda \bar K^\lambda_{\mu\nu}(t)\,s_\lambda
&&\text{(their ``generalized inverse Kostka''; the $s$ basis).}
\end{align}
Their $\bar K$ specialises at $\mu=0$ to the expansion of $P_\nu$ in the
\textbf{Schur} basis. That is the $P\to s$ matrix. The inverse of the $P\to m$ matrix
is $m\to P$. These are different matrices, and the word ``inverse Kostka'' in their
title refers to the former.

\paragraph{The measurement.} In the full text of \texttt{1603.01815} (6662 extracted
lines) the string \texttt{monomial symmetric} occurs \textbf{once}, at the eleventh
line of the introduction, in the remark that Hall--Littlewood functions interpolate to
monomial symmetric polynomials at $t=1$. Every other occurrence of ``monomial'' in the
paper means \emph{a monomial in the lattice-model variables} $x_i$ (``the coefficient
of a single monomial in $x$'', ``a staircase monomial''), never a basis element
$m_\mu$. \texttt{monomial basis} and \texttt{m\_\{} occur zero times.

In \texttt{1909.10720} (3232 lines) the same holds: \texttt{monomial symmetric} occurs
once, again as the $t=1$ interpolation remark, and the only occurrence of
\texttt{inverse Kostka} anywhere in the paper is \textbf{in the bibliography}, as the
title of the WZJ reference. That paper proves a Pieri rule and associativity for a
honeycomb formula for $f^\lambda_{\mu\nu}(t)$ --- again purely the $P$ basis.

\paragraph{Why I am willing to assert a null here.} I have been wrong before by
reporting a grep null that was an encoding artefact, so I ran the probe with controls.
Positive controls in \texttt{1603.01815}: \texttt{Hall} 55, \texttt{Kostka} 15,
\texttt{puzzle} 61, \texttt{Schur} 33. In \texttt{1909.10720}: \texttt{Hall} 43,
\texttt{honeycomb} 61. A planted nonsense string returned 0. The search terms here are
pure ASCII, so the en-dash class of false null does not apply. The nulls are real.

\section{So your $t=0$ edge is classical, not WZJ}

Since $P_\lambda(x;0)=s_\lambda$, at $t=0$ your object degenerates to
$[h_\mu]s_\lambda$, which is the Jacobi--Trudi / inverse-Kostka number --- classical.
The content of your law is therefore \emph{not} at $t=0$; it is the $(1-q)$-adic
valuation statement at general $q$, where the $t=0$ case is only the boundary value.
I would frame the novelty claim that way round.

\paragraph{Pointers I have \emph{not} verified at source.} Offered as leads, not
citations: Macdonald \emph{Symmetric Functions and Hall Polynomials} \S{}III.6 for the
$P\to m$ transition matrix itself; E\u{g}ecio\u{g}lu--Remmel (1990) for the combinatorial
inverse Kostka numbers via special rim-hook tabloids; and, closest to a
\emph{valuation} statement on Hall polynomials, Butler's \emph{Subgroup lattices and
symmetric functions}, which does work $p$-adically. If $\ell(\lambda)-\kappa(\lambda,\mu)$
is in the literature I would expect it there or in the cocharge
Kostka--Foulkes line (Lascoux--Sch\"utzenberger), not in the integrable-lattice line.
Treat all four as unchecked.

\section{Your question (b): the locator will not support the claim}

You propose that Theorem C cite \textbf{DFK \texttt{1505.01657} Cor.\ 5.18} for the
$t=0$ edge, ``whose $M_{k,1}$ is $E_k|_{t=0}$ literally''. That locator is
edition-confused, and I resolved this by label-probe compile on 2026-10-04:
\begin{itemize}
\item In the \textbf{arXiv} edition of \texttt{1505.01657}, the operator product
$\chi_n=\prod M\cdots\cdot 1$ --- the statement you want --- is \textbf{Cor.\ 5.8},
p.\ 20.
\item \textbf{Cor.\ 5.18}, p.\ 27, of that same edition is
$\chi_n=\lim P_\lambda^{q,t}$, and contains \textbf{no $M$ operator at all}.
\item The ``Corollary 18'' cited in \texttt{1908.00806} is the \textbf{journal}
numbering (\emph{Transform.\ Groups} \textbf{23}(2), 2018), which equals arXiv
Cor.\ 5.8. I believe that is where the number came from.
\end{itemize}
So: cite \textbf{Cor.\ 5.8} and name the edition, or cite ``Cor.\ 18 (journal
edition)''. As written, a referee checking the arXiv version finds a statement with no
$M$ in it. This is your own Risk 2 (``DFK15 normalization PENDING''); I consider it
closed by the above, in your favour --- the result you want \emph{is} in DFK, under a
different number.

\section{What I have not done}

I have not read \S5 of \texttt{1603.01815} (their rank-two fermion--boson model for
$\bar K$) in detail; only \S1--\S1.2, the definitions, and the closing normalisation.
If you want the $\bar K^\lambda_{\mu\nu}(t)$ puzzle formula itself quoted precisely,
ask and I will do that read properly. I have also not verified any of the four
pointers in \S2.

\end{document}
